\section{Limitations and open questions}
\label{sec:discussion-limitations-and-open-questions}

The risk comparisons leave several questions about joint changes in outcome arrival and alphabet size. The unrestricted upper bound in \cref{obj:thm:unrestricted-mixed-count-upper} provides an envelope across arrival regimes, while \cref{obj:thm:matched-minimax-frontier,obj:thm:unrestricted-fixed-q-minimax-frontier} give matching comparisons under their respective conditions. Matching the unrestricted envelope throughout intermediate moving-arrival and alphabet regimes remains open. This requires lower bounds that track simultaneous changes in the arrival floor and the number of baseline categories.

A related question concerns the distinction between a uniform near-complete requirement and a requirement for a particular alphabet sequence. \Cref{obj:prop:unrestricted-uniform-near-complete-elbow} characterizes the uniform comparison, and \cref{obj:thm:unrestricted-large-alphabet-deficit-lower} supplies the large-alphabet obstruction. These results leave open the precise near-complete requirement along individual smaller-alphabet sequences. A sequence-specific characterization would determine how the necessary degree of ascertainment depends on the growth of the alphabet.

For uncertainty quantification, \cref{obj:thm:connected-interval-frontier} compares uniformly honest connected intervals under the rare-arrival restriction in \cref{obj:ass:rare-arrival-slice}. Extending matching expected-length comparisons beyond that restriction is another open problem. Such an extension would require interval lower bounds appropriate to the additional arrival regimes, together with procedures attaining the corresponding lengths under uniform coverage.

The observational comparison in \cref{obj:thm:zeng-fixed-positivity-polynomial-frontier} concerns a fixed positivity margin strictly between zero and one half, as specified in \cref{obj:def:zeng-discrete-ate-class}. The endpoint at which every occupied category has treatment propensity one half calls for a separate risk characterization. Determining how this endpoint relates to the strict-interior comparisons would clarify the role of treatment-propensity variation in the observational problem.

The conditional arrival assumption in \cref{obj:ass:arrival-mar} also defines an important boundary of the model. Extending the analysis beyond missing at random requires specifying how arrival may depend on the outcome after conditioning on the recorded covariates, assignment, and surrogate. Alternative identifying restrictions or sensitivity parameters would then determine the target and the precision comparison. The present risk results do not supply those extensions.

Statistical attainment and practical evaluation are separate questions. The procedure in \cref{obj:def:mixed-count-estimator} averages \(3^{k_0}\) stream assignments at each retained prefix length \(k_0\), and the observational procedure in \cref{obj:thm:zeng-fixed-positivity-polynomial-frontier} additionally averages over \(2^n\) synthetic assignments. Their risk guarantees concern those exact procedures. Moreover, the polynomial branch begins only when \(\log(e+nq)\ge128\), or \(nq\ge e^{128}-e\). These scales make the displayed formulas mathematical attainment constructions rather than evidence of numerical usefulness at ordinary sample sizes. Computational evaluation should seek more efficient representations and quantify additional error from approximations. Prospective clinical evaluation must also justify the supplied arrival floor used for tuning and coverage, examine the choice of baseline categories and the conditional arrival assumption, and compare with established trial analyses.

\citet[abstract 300411]{Kennedy2019DiscreteAbstract} earlier announced average-causal-effect estimation with arbitrarily high-dimensional discrete confounders, including a variation with sample-size-dependent weakening of propensity-score positivity. The available abstract does not state its precise model class, rate, loss, or attaining construction. It therefore supports priority for studying that shrinking-positivity observational setting but cannot settle theorem-level overlap with the fixed-margin observational result here. A comparison with a complete theorem statement or manuscript remains open.
